% =====================================================================
%  Seance 1 -- B. Anatomie d'une infrastructure (frames 14--24)
% =====================================================================
\section{Anatomie d'une infrastructure et ses indicateurs}

\begin{frame}{Les couches d'une infrastructure moderne}
\begin{center}
\begin{tikzpicture}[node distance=2.5mm,
  layer/.style={rectangle,rounded corners=1.5mm,draw,align=center,
    font=\scriptsize,minimum height=7.5mm,text width=13.2cm}]
\node[layer,draw=byBlue,fill=byBlue!8]   (app)  {\textbf{Applications \& services} : API, bases de données, orchestrateurs --- ce que l'utilisateur consomme};
\node[layer,draw=byCyan,fill=byCyan!8,above=of app]  (cont) {\textbf{Conteneurs} : exécution isolée, limites CPU/RAM par service (Docker, Kubernetes)};
\node[layer,draw=byPurple,fill=byPurple!6,above=of cont] (virt) {\textbf{Virtualisation} : hyperviseur, machines virtuelles, allocation dynamique};
\node[layer,draw=byGreen,fill=byGreen!8,above=of virt] (mater){\textbf{Matériel} : CPU, RAM, disques, alimentation, refroidissement, salle};
\node[layer,draw=byOrange,fill=byOrange!8,below=of app] (reseau){\textbf{Réseau} : switches, pare-feu (\py{fw-edge}), liens opérateur, 4G/5G de NOSYS};
\end{tikzpicture}
\end{center}
\vspace{1mm}
\begin{itemize}
\item Chaque couche \textbf{émet ses propres données} : métriques matériel, télémétrie
de conteneurs, journaux applicatifs, flux réseau.
\item Un incident utilisateur est souvent la \textbf{signature} d'un problème deux couches
plus bas : c'est pourquoi on collecte tout, puis on apprend les corrélations.
\end{itemize}
\end{frame}

\begin{frame}{Métriques système : ce que mesurent les serveurs}
\begin{columns}[t]
\begin{column}{0.52\textwidth}
\begin{block}{Les cinq métriques du fichier NOSYS}
\begin{itemize}
\item \py{cpu\_pct} --- occupation processeur (\%).
\item \py{ram\_pct} --- occupation mémoire (\%).
\item \py{temperature\_c} --- température salle/machine (°C).
\item \py{conso\_kwh} --- consommation électrique (kWh par fenêtre).
\item \py{capacite\_cpu} --- capacité normalisée du serveur (0--1) :
sert au TP pour le \textbf{ratio d'utilisation effective}.
\end{itemize}
\end{block}
\begin{ebox}[Valeurs réelles du mois NOSYS]
CPU moyen du parc : \textbf{20,5 \%} (min 14,4 \%, max 24,6 \%) ;
RAM moyenne 41 à 61 \% ; température maximale \textbf{42,7 °C},
moyenne 32,5 °C. Source : \file{key\_numbers\_git1.json}.
\end{ebox}
\end{column}
\begin{column}{0.44\textwidth}
\begin{kbox}[Bonnes pratiques de collecte]
\begin{itemize}
\item Pas \textbf{constant} (ici 5 min) : indispensable pour les séries temporelles.
\item Horodatage en \textbf{UTC} côté collecte, conversion à l'affichage.
\item Ne jamais moyenner les \textbf{percentiles} : p99 de moyennes $\ne$ p99.
\item Un p98 de 42,7 °C n'est pas un souci ; une \textbf{dérive} de +0,3 °C/semaine si.
\end{itemize}
\end{kbox}
\begin{wbox}[Ce que le CPU seul ne dit pas]
Un serveur à 90 \% de CPU peut être sain (file d'attente vide)
et un serveur à 40 \% saturé (latence explosive).
\textbf{Toujours croiser} avec la latence et la perte.
\end{wbox}
\end{column}
\end{columns}
\end{frame}

\begin{frame}{Métriques réseau : latence, gigue, perte}
\begin{block}{Les trois mesures de base}
\begin{itemize}
\item \textbf{Latence} (ms) : temps de traversée. On suit la \textbf{distribution},
pas la moyenne : p50, p95, \textbf{p99} --- les utilisateurs mécontents sont dans la queue.
\item \textbf{Gigue} (jitter) : variation de la latence ; tue la voix/vidéo même
quand la latence moyenne est bonne.
\item \textbf{Perte de paquets} (\%) : chaque \% de perte coûte des retransmissions
$\Rightarrow$ effet multiplicatif sur la latence effective.
\end{itemize}
\end{block}
\begin{columns}[t]
\begin{column}{0.48\textwidth}
\begin{ebox}[NOSYS : chiffres réels]
Latence moyenne du parc : \textbf{12,2 ms} ;
\textbf{p99 = 16,3 ms} ; maximum observé \textbf{156,1 ms}
(pendant l'incident du 9/09, srv-01).
Perte : p99 = 0,17 \%. \textbf{Le SLO de 100 ms est tenu hors incidents.}
\end{ebox}
\end{column}
\begin{column}{0.48\textwidth}
\begin{kbox}[Moyenne vs p99 : l'exemple qui convainc]
La moyenne de 12 ms cache que 1 requête sur 100 dépasse 16 ms.
Sur 10 millions de requêtes/jour, cela fait \textbf{100 000 requêtes lentes} :
c'est là que vivent les plaintes clients --- et c'est pourquoi on supervise
le p99, jamais la moyenne seule.
\end{kbox}
\end{column}
\end{columns}
\end{frame}

\begin{frame}{SLI, SLO, SLA : ne plus jamais les confondre}
\begin{center}
\small
\begin{tabular}{@{}p{1.2cm}p{4.9cm}p{6.4cm}@{}}
\toprule
\textbf{Terme} & \textbf{Définition} & \textbf{Exemple NOSYS}\\
\midrule
\textbf{SLI} & L'\emph{indicateur} : ce qu'on mesure réellement. &
Latence p99 mesurée chaque minute ; taux de succès des requêtes.\\
\addlinespace[1pt]
\textbf{SLO} & L'\emph{objectif interne} posé sur le SLI. &
« 99,9 \% des requêtes $<$ 100 ms sur 30 jours glissants ».\\
\addlinespace[1pt]
\textbf{SLA} & Le \emph{contrat} avec le client, avec conséquences
(pénalités, crédits). & « Disponibilité 99,9 \% mensuelle, sinon crédit de 5 \% ».\\
\bottomrule
\end{tabular}
\end{center}
\vspace{1mm}
\begin{columns}[t]
\begin{column}{0.48\textwidth}
\begin{kbox}[La hiérarchie opérationnelle]
Le SLA \textbf{contraint} le SLO, le SLO \textbf{contraint} l'ingénierie.
En practice : on choisit le SLO avec une marge sous le SLA
--- c'est le « budget d'erreur » (error budget) que le ML de la
séance 4 aidera à défendre.
\end{kbox}
\end{column}
\begin{column}{0.48\textwidth}
\begin{qbox}[Mini-question]
NOSYS mesure une latence p99 de 16,3 ms. Le SLO est « p99 $<$ 100 ms ».
Quelle est la marge ?\\
(a) 6$\times$ \quad (b) 83,7 ms \quad (c) 612 \%
\end{qbox}
\end{column}
\end{columns}
\end{frame}

\begin{frame}{La disponibilité : ce qu'un « neuf » coûte vraiment}
\begin{center}
\includegraphics[width=0.88\textwidth]{f06_sla.png}
\end{center}
\begin{itemize}
\item Chaque 9 supplémentaire \textbf{divise le temps d'indisponibilité autorisé par 10} :
passer de 99,9 \% à 99,99 \% = passer de 8,8 h à 53 min d'arrêt \emph{annuel}.
\item Conséquence : à 99,999 \%, une panne de \textbf{30 min} consomme
5,7 \textbf{ans} de budget d'erreur. On ne « rattrape » pas une panne à ce niveau.
\end{itemize}
\end{frame}

\begin{frame}{Calculs de disponibilité : le rituel}
\begin{columns}[t]
\begin{column}{0.48\textwidth}
\begin{block}{Les deux formules}
\[
\text{Disponibilité} = \frac{\text{temps de service correct}}{\text{temps total}}
= 1 - \frac{\text{indisponibilité}}{\text{temps total}}
\]
\[
\text{MTBF} = \frac{\text{temps total de service}}{\text{nombre de pannes}}
\qquad
\text{MTTR} = \frac{\text{temps cumulé de réparation}}{\text{nombre de pannes}}
\]
\end{block}
\begin{ebox}[Exemple calculé sur NOSYS]
16 incidents $\times$ 8 serveurs = 16 blocs sur le mois, durées
35 et 50 min (0,58 h et 0,83 h). Indisponibilité cumulée du parc :
$8 \times (0{,}58 + 0{,}83) = 11{,}3$ h sur $8 \times 720$ h
$\Rightarrow$ disponibilité \textbf{99,80 \%}.
\end{ebox}
\end{column}
\begin{column}{0.48\textwidth}
\begin{wbox}[Conclusion à retenir]
Avec 2 incidents/serveur/mois, NOSYS \textbf{rate son SLA de 99,9 \%}
(99,80 \% atteint). Toute la mission ML des séances 3--4 se résume à :
\textbf{transformer ces 11,3 h de panne réactive en prévention}.
\end{wbox}
\begin{qbox}[Pour le prochain cours]
Si un prédictif détecte 70 \% des incidents 30 min à l'avance et que
l'équipe les évite, quelle devient la disponibilité ?
(faites le calcul --- réponse en séance 4.)
\end{qbox}
\end{column}
\end{columns}
\end{frame}

\begin{frame}{Quiz 2 --- les indicateurs (5 min)}
\begin{qbox}[3 questions]
\begin{enumerate}
\item Le p99 de latence est de 16 ms. Combien de requêtes sur 1 000 dépassent ce seuil ?\\
(a) 1 \quad (b) 10 \quad (c) 99
\item Un serveur a 6 pannes de 20 min sur l'année (8 760 h). MTTR = 20 min ;
MTBF $\approx$ ?\\
(a) 1 460 h \quad (b) 8 760 h \quad (c) 120 min
\item Pourquoi suivre le p99 plutôt que la moyenne ?\\
(a) le p99 est plus facile à calculer \quad (b) la moyenne masque la queue
de distribution où se trouvent les utilisateurs pénalisés \quad
(c) le p99 est exigé par le SLA uniquement
\end{enumerate}
\end{qbox}
\vspace{1mm}
\begin{solbox}[Réponses]
\textbf{1 : b} --- p99 = 99 \% en dessous, 1 \% au-dessus = 10 sur 1 000.
\quad
\textbf{2 : a} --- MTBF = 8 760 / 6 $\approx$ 1 460 h.
\quad
\textbf{3 : b} --- la moyenne est dominée par la masse de requêtes rapides.
\end{solbox}
\end{frame}

\begin{frame}{À retenir de la partie anatomie}
\begin{kbox}[Les 5 idées de cette section]
\begin{enumerate}
\item Une infrastructure est une \textbf{pile de couches} ; chaque couche émet
ses données --- et un symptôme applicatif a souvent sa cause deux couches plus bas.
\item Superviser la \textbf{distribution} des métriques (p99), pas la moyenne.
\item \textbf{SLI} = ce qu'on mesure ; \textbf{SLO} = l'objectif interne ;
\textbf{SLA} = le contrat. Ne jamais les confondre devant un client.
\item La disponibilité se \textbf{calcule} : 99,9 \% = 8,8 h/an ;
chaque 9 supplémentaire divise le budget par 10.
\item NOSYS est à \textbf{99,80 \%} : l'UE entière vise à remonter ce chiffre
par la prédiction et la prévention.
\end{enumerate}
\end{kbox}
\end{frame}
